% ---
% id: receiver/lem:fock-block-spectrum
% kind: lem
% title: "Spectrum of a lifted no-click operator on a photon-number block"
% summary: >-
%   On the n-photon block, the lift of 1 − A has as eigenvalues the products of
%   n factors 1 − λ over the eigenvalues λ of A, so it lies between the n-th
%   powers of 1 − λmax and 1 − λmin, both attained. The analytic high-photon
%   bound of the receiver starts from these two extremes.
% prominence: standard
% status: implemented
% uses: [lem:fock-functor]
% mentions: [thm:port-grams]
% assumes: []
% refs: []
% provenance: newly-derived
% implements: certiqs_key.receiver.fock:block_spectrum_from_eigenvalues
% lean:
% ---

For $A$ the sum of the \tref[port Gram matrices]{port Grams} of a set of ports,
$\Fock(\id-A)$ is the operator for the event that no port of the set clicks,
pulled back to the declared input modes (Theorem~\ref{thm:port-grams}\,(iii));
hence the name lifted no-click operator.

\begin{lemma}
\label{lem:fock-block-spectrum}
Let $A$ be Hermitian on $\mathcal{K}_0$ with $0\preceq A\preceq\id$ and
eigenvalues $\lambda_1,\dots,\lambda_d$ ($d=\dim\mathcal{K}_0$). Then
$\Fock(\id-A)$ is a positive semidefinite contraction, its $n$-photon block
$\Fock(\id-A)^{(n)}=\big((\id-A)^{\otimes n}\big)\big|_{\Symn{n}}$ has spectrum
\begin{equation}\label{eq:spectrum}
  \Big\{\textstyle\prod_{j=1}^{d}(1-\lambda_j)^{m_j}\ :\
  \vec m\in\mathbb{N}_0^{d},\ |\vec m|=n\Big\},
\end{equation}
and consequently
\begin{equation}\label{eq:blockbounds}
  \big(1-\lmax(A)\big)^{n}\,\id\ \preceq\ \Fock(\id-A)^{(n)}\ \preceq\
  \big(1-\lmin(A)\big)^{n}\,\id ,
\end{equation}
both bounds being attained as eigenvalues.
\end{lemma}

\begin{proof}
Diagonalize $A=U^\dagger\Lambda U$ with $U$ unitary and
$\Lambda=\operatorname{diag}(\lambda_1,\dots,\lambda_d)$; then
$\id-A=U^\dagger(\id-\Lambda)U$ and Lemma~\ref{lem:fock-functor}\,(ii),(iii) give
$\Fock(\id-A)=\Fock(U)^\dagger\Fock(\id-\Lambda)\Fock(U)$ with $\Fock(U)$
unitary and photon-number preserving, so the two operators have the same
spectrum block by block. In the occupation basis built from the eigenbasis of
$A$, the product vector $\ket{e_{j_1}}\otimes\cdots\otimes\ket{e_{j_n}}$ is an
eigenvector of $(\id-\Lambda)^{\otimes n}$ with eigenvalue
$\prod_r(1-\lambda_{j_r})$, and that eigenvalue is unchanged by permuting the
factors; hence the symmetrized basis vector $\ket{\vec m}$ is an eigenvector of
$(\id-\Lambda)^{\otimes n}|_{\Symn{n}}$ with eigenvalue
$\prod_j(1-\lambda_j)^{m_j}$. These $\ket{\vec m}$, $|\vec m|=n$, form an
orthonormal basis of $\Symn{n}$, so \eqref{eq:spectrum} is the full spectrum.
Each factor satisfies $1-\lambda_j\in[1-\lmax(A),1-\lmin(A)]\subseteq[0,1]$,
so a product of $n$ of them lies in
$[(1-\lmax(A))^n,(1-\lmin(A))^n]$, which is \eqref{eq:blockbounds}; the
endpoints are attained by the occupation vectors putting all $n$ photons on an
eigenvector for $\lmax(A)$, respectively for $\lmin(A)$. Positivity and
contractivity follow from \eqref{eq:blockbounds} and
Lemma~\ref{lem:fock-functor}\,(iii).
\end{proof}
